\scanpage{40}
De plus, (1) (2) (3) [\unsure{impliquant} l'existence d'un $G$ simple sur $S$ donnant $\mathfrak g$] prouvent aussi [pour $G$ connexe] que $G$ n'a pas de variations infinitésimales non triviales.

Question du relèvement de la structure $G$ \inserted{(connexe)}, équivaut à la question du relèvement de la structure $\mathfrak g$. On trouve une obstruction dans $H^3(\mathfrak g,\mathfrak g)$. Donc

\deleted{$H^3$}

(3) Si $H^3(\mathfrak g,\mathfrak g)=0$, alors on peut relever de façon essentiellement unique $\mathfrak g$, ou $G$ de façon équivalente.

\uline{Question} Définir
\[
H^{\unsure{*}}(G,\mathfrak g)\longrightarrow H^*(\mathfrak g,\mathfrak g)
\]
[plus généralement, si on a une représentation linéaire de $G$\ldots]. Suite spectrale reliant $H^*(G,\ )$ et $H^*(\mathfrak g,-)$ ?? Elle exprimerait la ``différence'' entre le global, et l'infinitésimal.

Supposons $G$ semi-simple adjoint. A-t-on alors
\[
H^*(G,\mathfrak g)=H^*(\mathfrak g,\mathfrak g)=0\qquad??
\]
\margin{39}
